\begin{titlepage}
\centering
{\Large Notas didácticas de revisión técnica\par}
\vspace{1.0cm}
{\huge\bfseries Repaso de los informes técnicos de Agent: Kimi K2, ChatGPT Agent, Qwen3-Coder y Manus}\par
\vspace{0.5cm}
{\Large Zhang Xiaojun conversa con Zheng Boyuan (郑博元)\par}
\vspace{0.35cm}
{\large Elaborado a partir del video público del Zhang Xiaojun Podcast, una transcripción hecha en GPU, los informes técnicos y entradas de blog oficiales, y diagramas conceptuales propios\par}
\vspace{0.3cm}
{\large 2025-07-30\par}
\vspace{0.85cm}
\includegraphics[width=0.88\textwidth,height=0.42\textheight,keepaspectratio]{cover.jpg}\par
\vfill
\begin{tcolorbox}[width=0.92\textwidth, colback=black!2!white, colframe=black!55, sharp corners]
\textbf{Título del video}: 110. Explicación sección por sección del informe de Kimi K2, comparado con ChatGPT Agent, Qwen3-Coder y otros: «El poder de la ingeniería de sistemas»\par
\textbf{Canal del video}: Zhang Xiaojun Podcast\par
\textbf{Duración del video}: 02:20:45\par
\textbf{Enlace del video}: \href{\videourl}{\nolinkurl{\videourl}}\par
\textbf{Descripción de los materiales}: El video no tiene subtítulos disponibles en YouTube; el cuerpo de la nota se elaboró a partir de la transcripción de faster-whisper large-v3 con CUDA, los capítulos del video, la descripción de YouTube, el informe técnico de Kimi K2, la entrada de blog de Qwen3-Coder, la entrada de blog de Manus sobre ingeniería de contexto y diagramas didácticos propios. La imagen fija de portada del video no se repite en el cuerpo de la nota.\par
\textbf{Página del proyecto}: \href{\repourl}{\nolinkurl{\repourl}}\par
\end{tcolorbox}
\end{titlepage}

\tableofcontents
\newpage

